\section{A common first-moment framework}
\label{sec:framework}

This section isolates the calculation shared by all SPIN families. We first
define the complete sampled encoder. We then count low-weight outputs of
nonzero messages.

Fix lengths $K,N\ge1$. The construction samples an outer encoder
$O:\F_2^K\to\F_2^N$ and a pair $(\Pi,I)$, where $\Pi$ is a permutation of
$[N]$ and $I:\F_2^N\to\F_2^N$ is linear. Assume that $O$ is independent
of $(\Pi,I)$. The realized SPIN encoder is
\[
  E:=I\circ\Pi\circ O.
\]

Fix a finite class set $\mathcal T$ and a class map
$c:\F_2^N\to\mathcal T$. The class map records the information about an
outer word that remains relevant after the outer stage.

For a fixed outer word $u\in\F_2^N$, define its conditional failure
probability
\begin{equation}
  q_d(u)
  :=
  \Prob_{\Pi,I}[\wt(I(\Pi(u)))\le d].
  \label{eq:conditional-inner-failure}
\end{equation}
The probability is over the interleaver and inner setup; the word $u$ is
fixed. For every class $\tau\in\mathcal T$, let $Q_\tau(d)$ be any certified
envelope satisfying
\begin{equation}
  q_d(u)\le Q_\tau(d)
  \qquad\text{for every $u$ with $c(u)=\tau$}.
  \label{eq:class-failure-envelope}
\end{equation}
Set $Q_\tau(d):=0$ when the class has no preimage. The class need not determine
the exact failure probability. It need only expose enough information to prove
the uniform bound in~\eqref{eq:class-failure-envelope}.

For $\tau\in\mathcal T$, define the expected outer multiplicity
\begin{equation}
  A_\tau^{\mathrm{out}}
  :=
  \E_O\!\left[
    \sum_{x\in\F_2^K\setminus\{0\}}
    \1\{c(O(x))=\tau\}
  \right].
  \label{eq:outer-class-multiplicity}
\end{equation}

\begin{theorem}[Class-indexed first-moment envelope]
\label{thm:class-first-moment}
For $d\in\{0,\ldots,N\}$, define
\[
  Z_d
  :=
  \bigl|\{x\in\F_2^K\setminus\{0\}:
  \wt(E(x))\le d\}\bigr|.
\]
Then
\begin{equation}
\begin{split}
  \Prob[E\text{ is not injective or }d_{\min}(E(\F_2^K))\le d]
  \le
  \E[Z_d]
  &=
  \E_O\!\left[
    \sum_{x\in\F_2^K\setminus\{0\}}q_d(O(x))
  \right]\\
  &\le
  \sum_{\tau\in\mathcal T}
  A_\tau^{\mathrm{out}}Q_\tau(d).
\end{split}
  \label{eq:class-first-moment}
\end{equation}
If $q_d$ is constant on every class and $Q_\tau(d)$ equals that constant,
then the final inequality is an equality.
\end{theorem}

\begin{proof}
The bad event in~\eqref{eq:class-first-moment} occurs exactly when
$Z_d\ge1$. Indeed, a nonzero kernel message contributes output weight zero,
and every nonzero image word of weight at most $d$ has a nonzero preimage.
Markov's inequality gives $\Prob[Z_d\ge1]\le\E[Z_d]$.

For the identity, expand $Z_d$ as a sum over nonzero messages and condition
on the outer realization. Independence of $O$ from $(\Pi,I)$ gives
\[
  \E[Z_d]
  =
  \E_O\!\left[
    \sum_{x\ne0}q_d(O(x))
  \right].
\]
Group the terms by $\tau=c(O(x))$ and apply
\eqref{eq:class-failure-envelope}. Equation
\eqref{eq:outer-class-multiplicity} gives the final inequality. Equality holds
under the condition in the theorem statement.
\end{proof}

If $\mathcal G$ is an event determined only by the outer setup, the theorem
also applies under the conditional law $O\mid\mathcal G$. The pair $(\Pi,I)$
remains independent of the conditioned outer encoder. This observation lets a
proof first select an outer spectrum event and then certify route envelopes on
that event.

\subsection{Uniform-interleaver specialization}

Suppose $\Pi$ is uniform over all permutations of $[N]$. If $u$ has
weight $w$, then $\Pi(u)$ is uniform over the Hamming slice
\[
  \mathcal H_{N,w}:=\{v\in\F_2^N:\wt(v)=w\}.
\]
Indeed, every vector in $\mathcal H_{N,w}$ has exactly
$w!(N-w)!$ preimages among the permutations of $u$.

Let $U_w$ be uniform over $\mathcal H_{N,w}$, independent of the inner
setup, and define
\begin{equation}
  p_w^{\mathrm{in}}(d)
  :=
  \Prob_{U_w,I}[\wt(I(U_w))\le d].
  \label{eq:inner-slice-tail}
\end{equation}
Also define
\[
  A_w^{\mathrm{out}}
  :=
  \E_O\bigl[
    |\{x\ne0:\wt(O(x))=w\}|
  \bigr].
\]

\begin{corollary}[Weight-indexed first moment]
\label{cor:weight-first-moment}
Under a uniform interleaver independent of $O$ and $I$,
\[
  \E[Z_d]
  =
  \sum_{w=0}^N
  A_w^{\mathrm{out}}p_w^{\mathrm{in}}(d).
\]
\end{corollary}

Accumulator SPIN and Random SPIN use
Corollary~\ref{cor:weight-first-moment}. Structured SPIN uses
Theorem~\ref{thm:class-first-moment} with occupation and row-weight classes.
The structured-route analysis supplies a uniform failure envelope for every
such class; it does not require the conditional probability to be identical
for all words in the class.
